\section{Hiện thực hệ thống}

\subsection{Khối 1: Hiện thực LED + LCD, Ngoại vi (Triển khai trên Yolo UNO 2)}
Khối chức năng ngoại vi cảnh báo và hiển thị cục bộ được hiện thực hóa toàn bộ trên bo mạch \textbf{Yolo UNO Khối 2 (Display Node)}. Toàn bộ các chân GPIO kích hoạt LED, NeoPixel và bus I2C của LCD được cấu hình thông qua PlatformIO. Mã nguồn được cấu trúc lại hoàn toàn để xóa bỏ 100\% biến toàn cục (global variables) thông qua việc đóng gói dữ liệu vào cấu trúc dữ liệu con trỏ \texttt{SystemData} hoặc truyền an toàn qua hàng đợi/semaphore.

\subsubsection{Hiện thực tác vụ LED đơn (Task 1)}
Tác vụ quản lý đèn LED đơn chỉ thị nhiệt độ được khởi chạy với một vòng lặp vô hạn \texttt{for(;;)}. Tác vụ liên tục rơi vào trạng thái chặn (Blocked State) để tiết kiệm tài nguyên CPU thông qua hàm FreeRTOS API:
\begin{lstlisting}[language=C++]
xSemaphoreTake(xLedSemaphore, portMAX_DELAY);
\end{lstlisting}
Ngay khi tác vụ nhận truyền thông giải phóng semaphore, Task 1 sẽ thức dậy, đối chiếu giá trị nhiệt độ trích xuất từ gói tin và hiện thực thuật toán cấu hình chu kỳ nhấp nháy (\texttt{vTaskDelay}) qua 3 chế độ dải đo rõ ràng:
\begin{itemize}
    \item \textbf{Chế độ an toàn ($T < 30^{\circ}$C):} LED tắt hoàn toàn hoặc nhấp nháy với chu kỳ rất dài (bật 100ms, tắt 2000ms).
    \item \textbf{Chế độ cảnh báo ($30^{\circ}\text{C} \le T \le 40^{\circ}$C):} LED nhấp nháy chu kỳ trung bình (bật 500ms, tắt 500ms).
    \item \textbf{Chế độ nguy cấp ($T > 40^{\circ}$C):} LED băm xung hoặc nhấp nháy tần số cao (bật 100ms, tắt 100ms) để biểu thị trạng thái khẩn cấp.
\end{itemize}

\subsubsection{Hiện thực tác vụ điều khiển NeoPixel (Task 2)}
Thay vì sử dụng Task Notifications như các hệ thống cũ, kỹ thuật đồng bộ hóa của dải LED RGB NeoPixel (WS2812B) được hiện thực nhất quán bằng Semaphore tương tự như LED đơn:
\begin{lstlisting}[language=C++]
xSemaphoreTake(xNeoSemaphore, portMAX_DELAY);
\end{lstlisting}
Khi lấy được token, hệ thống gọi hàm thư viện \texttt{pixels.setPixelColor()} để ánh xạ trực tiếp mức độ ẩm sang 3 mẫu màu sắc định trị cụ thể: màu Đỏ (độ ẩm thấp, không khí khô), màu Xanh lá (độ ẩm lý tưởng), và màu Xanh dương (độ ẩm quá cao, môi trường ẩm ướt).

\subsubsection{Hiện thực tác vụ hiển thị LCD I2C (Task 3)}
Tác vụ LCD hiện thực cơ chế Mutex nhằm cô lập hoàn toàn bus I2C, ngăn hiện tượng tranh chấp xung nhịp (Race Condition) gây nhiễu chữ khi màn hình ghi dữ liệu. Tác vụ giải mã chuỗi dữ liệu nhận được và chuyển đổi linh hoạt 3 giao diện hiển thị: dòng chữ ``STATUS: NORMAL'' trên nền thông số thực tế, dòng ``STATUS: WARNING'' đi kèm nháy đèn nền LCD, và dòng khóa màn hình ``STATUS: CRITICAL'' khi TinyML hoặc cảm biến kích hoạt ngưỡng nguy hại.

\begin{figure}[H]
    \centering
    % \includegraphics[width=0.85\textwidth]{fritzing_node2_wiring.png}
    \caption{Sơ đồ mạch lắp ráp Fritzing thực tế của các ngoại vi kết nối với bo mạch Yolo UNO 2}
    \label{fig:fritzing_node2}
\end{figure}

\subsection{Khối 2: Webserver (Triển khai trên Yolo UNO 2)}
Mã nguồn máy chủ Web Server phục vụ việc quản trị và tương tác điều khiển thiết bị cục bộ tại biên được hiện thực toàn bộ trên bộ xử lý của bo mạch \textbf{Yolo UNO Khối 2}.
\begin{itemize}
    \item \textbf{Cấu hình mạng AP:} Khởi tạo lớp kết nối không dây ở chế độ Điểm truy cập thông qua cú pháp \texttt{WiFi.softAP(ssid, password)}, đặt cấu hình dải IP tĩnh là \texttt{192.168.4.1}.
    \item \textbf{Hiện thực giao diện trực quan:} Toàn bộ file giao diện (\texttt{index.html}, \texttt{script.js}, \texttt{style.css}) được lưu giữ trực tiếp trong phân vùng bộ nhớ Flash (nhờ LittleFS hoặc PROGMEM) nhằm tối ưu RAM.
    \item \textbf{Xử lý Endpoint qua WebSocket:} Điểm khác biệt cốt lõi là hệ thống loại bỏ các phương thức HTTP GET/POST rườm rà, thay vào đó hiện thực hóa giao thức WebSocket để truyền tải luồng JSON tốc độ cao:
\begin{lstlisting}[language=C++]
AsyncWebSocket ws("/ws");
ws.onEvent(onEvent);
server.addHandler(&ws);
\end{lstlisting}
Khi người dùng tương tác trên Dashboard (thêm/xóa/bật/tắt Relay động), trình duyệt sẽ gửi gói tin JSON qua kênh WebSocket. Trình xử lý \texttt{onEvent} trên vi điều khiển lập tức phân tích cú pháp (Parse JSON) để điều khiển các cổng GPIO tương ứng theo thời gian thực mà không phát sinh độ trễ.
\end{itemize}

\subsection{Khối 3: CoreIOT + MQTT (Triển khai trên Yolo UNO 1)}
Để tách biệt luồng xử lý mạng diện rộng ra khỏi luồng điều khiển WebUI, khối tương tác đám mây được mã hóa và hiện thực độc quyền trên nền tảng \textbf{Yolo UNO Khối 1 (Sensor Node)}.
\begin{itemize}
    \item \textbf{Cấu hình mạng STA:} Kích hoạt Wi-Fi hoạt động ở chế độ Trạm bằng lệnh \texttt{WiFi.begin()} nhằm kết nối vào bộ định tuyến và lấy địa chỉ IP kết nối mạng Internet.
    \item \textbf{Hiện thực đẩy dữ liệu lên Cloud (Task 6):} Tác vụ định kỳ đóng gói giá trị số thực của nhiệt độ và độ ẩm thu thập được từ cảm biến vật lý DHT20 vào cấu trúc JSON tiêu chuẩn. Mã nguồn gọi thư viện ứng dụng ThingsBoard MQTT để thực thi phương thức xuất bản (publish) lên máy chủ đám mây CoreIOT (\url{https://app.coreiot.io/}). Gói tin gửi đi luôn được đính kèm \texttt{Device Token} để vượt qua hệ thống tường lửa bảo mật của máy chủ.
\end{itemize}

\subsection{Khối 4: TinyML (Triển khai trên Yolo UNO 1)}
Thuật toán Trí tuệ nhân tạo biên đòi hỏi công suất tính toán ma trận lớn, do đó được hiện thực cấu hình chạy độc lập trên \textbf{Yolo UNO Khối 1}, giải phóng hoàn toàn bộ nhớ cho các tác vụ hiển thị của bo mạch 2.

\subsubsection{Hiện thực TinyML Task}
Tác vụ TinyML (Task 5) được nhóm hiện thực hóa ghim cố định trên lõi CPU chuyên dụng (Core 1) của dòng chip lõi kép ESP32-S3 bằng hàm API mở rộng của FreeRTOS:
\begin{lstlisting}[language=C++]
xTaskCreatePinnedToCore(TinyML_Task, "TinyML_Task", 8192, NULL, 3, &xTinyMLTaskHandle, 1);
\end{lstlisting}
Điều này đảm bảo tiến trình AI hoạt động with vùng Stack cô lập 8KB mà không gây tràn bộ nhớ hay làm chậm chu kỳ lập lịch đọc cảm biến ở lõi 0.

\subsubsection{Hiện thực TensorFlow Lite Micro và Tensor Arena}
Hệ thống nạp khung sườn thư viện thông qua tệp tiêu đề bộ thông dịch \texttt{\#include "tensorflow/lite/micro/micro\_interpreter.h"}. Để triệt tiêu lỗi sập nguồn do phân mảnh RAM động (Heap Fragmentation), bộ nhớ Tensor Arena được hiện thực dưới dạng mảng byte tĩnh 8KB, căn chỉnh biên độ phần cứng 16-bit nghiêm ngặt:
\begin{lstlisting}[language=C++]
constexpr int kTensorArenaSize = 8 * 1024; // 8KB
alignas(16) static uint8_t tensor_arena[kTensorArenaSize];
\end{lstlisting}

\subsubsection{Nạp dữ liệu trực tiếp và Kịch bản kiểm thử tĩnh (Startup Test)}
Điểm nhấn khác biệt của thuật toán là việc loại bỏ hoàn toàn mã chuẩn hóa cửa sổ trượt phức tạp. Dữ liệu từ cảm biến DHT20 được ép kiểu và nạp thẳng vào con trỏ tensor dưới định dạng số thực \texttt{float32}:
\begin{lstlisting}[language=C++]
model_input->data.f[0] = current_temp; // Truyen thang Float32
model_input->data.f[1] = current_humi;
\end{lstlisting}
Đặc biệt, hệ thống hiện thực một vòng lặp \texttt{evaluate\_accuracy()} đánh giá tự động dựa trên 50 mẫu Test Dataset ngay khi khởi động. Thuật toán tự đối chiếu nhãn dự đoán với nhãn gốc (Ground Truth) để tính toán \textit{True Positives/False Positives}, từ đó xuất ra các hệ số Accuracy, Precision, Recall qua cổng nối tiếp trước khi vào chế độ vận hành chính thức.

\subsubsection{Hiện thực cơ chế cảnh báo và Đo lường độ trễ}
Sau khi gọi hàm \texttt{interpreter->Invoke()}, mã nguồn truy cập trực tiếp mảng tensor đầu ra để lấy hệ số điểm bất thường. Hệ thống hiện thực bộ lọc logic: nếu điểm bất thường (Anomaly Score) vượt ngưỡng cấu hình an toàn ($> 0.85$), hệ thống sẽ lập tức gắn cờ khẩn cấp vào gói tin truyền thông nội bộ. 
Đồng thời, hàm \texttt{esp\_timer\_get\_time()} được sử dụng để đo đạc thời gian thực thi suy luận (Inference Latency), hỗ trợ đánh giá hiệu suất mô hình chạy trên ESP32-S3.

\subsection{Khối 5: Giao tiếp MacAddress (Xác thực phân tán mạng 2 Board)}
Để hiện thực hóa cơ chế bảo mật định danh, chống các thiết bị giả mạo can thiệp vào mạng phân tán của hai board Yolo UNO, hệ thống hiện thực module đọc địa chỉ vật lý toàn cầu trực tiếp từ eFuse phần cứng bằng hàm API Espressif native:
\begin{lstlisting}[language=C++]
uint8_t mac_hardware[6];
esp_efuse_mac_get_default(mac_hardware);
\end{lstlisting}
Mã nguồn thực hiện định dạng mảng byte địa chỉ MAC thu được thành chuỗi ký tự Hex tiêu chuẩn (ví dụ: \texttt{"32:B9:E6:7A:8C:1F"}). 

Khi Yolo UNO Khối 1 đóng gói dữ liệu cảm biến DHT20 và kết quả TinyML gửi không dây sang Yolo UNO Khối 2, chuỗi MacAddress bắt buộc phải đính kèm vào phần tiêu đề gói tin (Packet Header). Tại Yolo UNO Khối 2, mã nguồn hiện thực một bộ lọc whitelist so sánh: nếu chuỗi MAC gửi tới không trùng khớp chính xác với địa chỉ vật lý của Yolo UNO 1 đã cấu hình cứng (hardcoded) từ trước, tác vụ nhận dữ liệu sẽ lập tức loại bỏ gói tin, triệt tiêu hoàn toàn nguy cơ giả mạo lệnh điều khiển trên mạng nội bộ.